\subsection{Uniformly finite coefficient diagrams}

\begin{proposition}[Proved: bounded finite coefficients]
\label{A:uniform-finite-coefficients}
Let $A$ be a nonempty connected small quasi-filtered category, and
let $M:A\to\Ab$ satisfy $|M(a)|\leq N$ for every object $a$, for
one finite integer $N$. Then $T_1^A(M)=0$. If $A$ is total, then
$T_i^A(M)=0$ for every $i>0$.

Consequently, a prime-field column module $Q_a$ with nonzero first
derived colimit over a quasi-filtered category must have unbounded
dimensions as the object varies.
\end{proposition}
\begin{proof}
Define $B$ to have the same objects as $A$ and the typed hom-sets
\[
 B(a,b)=\operatorname{im}\bigl(A(a,b)\longrightarrow
              \Hom_{\Ab}(M(a),M(b))\bigr).
\]
Composition is composition of homomorphisms. The quotient functor
$u:A\to B$ is the identity on objects and surjective on every
hom-set, and $M=u^*\overline M$ for a functor
$\overline M:B\to\Ab$. Conditions I and IV descend along $u$:
lift an ordinary span or parallel pair, solve the corresponding
problem in $A$, and apply $u$. Furthermore
\[
                   |B(a,b)|\leq N^{|M(a)|}\leq N^N.
\]
Thus \citet[Proposition~4.5]{IsbellMitchell1976} supplies an ordinary
post-equalizer for every finite parallel family in $B$. Connectedness
and I give common successors, so $B$ is filtered.

For $b\in B$, the comma category $b\mathbin\downarrow u$ is the
category of elements of the covariant $A$-set $B(b,u(-))$.
Its object $(b,\id_b)$ is weakly initial, by hom-surjectivity of
$u$. Hence this comma category is nonempty and connected. It is
simply connected under quasi-filteredness, by the audited
element-category form of I and IV; under totality it is weakly
contractible.

For completeness, the homological comparison uses these comma
homology groups, not ordinary finality alone. Choose the projective
right bar resolution $P_*\to\Delta\Z$ on $A$, with
\[
 P_n=\bigoplus_{a_0\to\cdots\to a_n}\Z[A(-,a_0)],
 \qquad K_* = \operatorname{Lan}_{u^{\mathrm{op}}}P_*.
\]
Left Kan extension sends each right representable to the
corresponding right representable on $B$. Thus $K_*$ is a
nonnegative complex of projective right $\Z B$-modules. At $b$ its
degree-$n$ group is
\[
 K_n(b)=\bigoplus_{a_0\to\cdots\to a_n}\Z[B(b,u(a_0))],
\]
with the nerve differential of $b\mathbin\downarrow u$.
Consequently $H_0K_*\cong\Delta\Z$ naturally and $H_1K_*=0$;
under totality, $H_jK_*=0$ for every $j>0$.

Coend adjunction identifies the two complexes
\[
 K_*\otimes_{\Z B}\overline M
       \cong P_*\otimes_{\Z A}u^*\overline M.
\]
The first-quadrant hyper-Tor spectral sequence is therefore
\[
 E^2_{p,q}=\Tor_p^{\Z B}(H_qK_*,\overline M)
       \ \Longrightarrow\ T_{p+q}^A(M).
\]
Since $B$ is filtered, $\Delta\Z$ is flat on $B$. Both terms of
total degree one vanish: $(p,q)=(1,0)$ by flatness and $(0,1)$ by
$H_1K_*=0$. This proves the first assertion. Under totality the
spectral sequence has only its $q=0$ row, and flatness proves all
positive-degree vanishings. Uniform finite dimension over $\F_p$
implies the required uniform cardinal bound, proving the consequence.
\end{proof}

This does not cover arbitrary finite-valued diagrams with unbounded
cardinalities, or arbitrary prime-field modules. In particular, a
finite set of coefficient vectors need not generate a uniformly
bounded subdiagram. Passing to a category of elements that pins a
tuple also need not preserve the relations of an original bar cycle;
Proposition~\ref{CD:pinning-counterexample} gives the relevant
warning about such a change of category.
